\begin{formulaBox}{Streuungsmaße}{ML I}
\begin{multicols}{2}

\subsection*{Standardabweichung}
$$s=\sqrt{s^2}$$

\subsection*{Varianz}

$$s^2=\frac{1}{n}\sum\limits_{i=1}^{n}(x_i-\bar{x})^2$$

\subsection*{Schätzung für Populationsvarianz}

$$\hat{\sigma}^2=\frac{1}{n-1}\sum\limits_{i=1}^n(x_i-\bar{x})^2=\frac{n}{n-1}\cdot s^2$$

Hinweis: Nach APA~7 steht $s$ für die mit $n-1$ berechnete Stichprobenstandardabweichung und $S$ für die deskriptive Standardabweichung mit $n$. Lehrbücher verwenden $s$ jedoch nicht einheitlich. Achten Sie daher stets auf die Definition im jeweiligen Kontext; in dieser Formelsammlung bezeichnet $s$ die oben mit $n$ definierte Größe. Die zugehörige Varianzformel mit $n-1$ ist erwartungstreu für die Populationsvarianz; für die gesamte Population wird mit $n$ gerechnet.

\subsection*{Spannweite (Range)} 

\begin{equation*}
R=x_\mathrm{max}-x_\mathrm{min}
\end{equation*}

\subsection*{Mittlerer absoluter Abstand}

$$\mathrm{MAD} = \frac{1}{n}\sum_{i=1}^n|x_i-\bar{x}|$$

Gemeint ist der \emph{Mittelwert} der absoluten Abstände, nicht die ebenfalls mit MAD abgekürzte \emph{median absolute deviation}.

\subsection*{Interquartilsabstand}

\begin{equation*}
\mathrm{IQA}=Q_{0.75}-Q_{0.25}
\end{equation*}
\vspace{0.1mm}

%\subsection*{Standardfehler}

%$$\hat{\sigma}=\sqrt{\hat{\sigma}^2}=\sqrt{\frac{n}{n-1}}\cdot s$$ 
\end{multicols}
\end{formulaBox}

\begin{formulaBox}{Summenalgebra}{ML I}
\begin{multicols}{2}
\[
\sum_{i=1}^{n}(a_i+b_i)
  = \sum_{i=1}^{n}a_i + \sum_{i=1}^{n}b_i
\]

\columnbreak

\[
\sum_{i=1}^{n}c\,a_i
  = c\sum_{i=1}^{n}a_i
\qquad
\sum_{i=1}^{n}c = n c
\]
\end{multicols}
\end{formulaBox}

\begin{formulaBox}{Transformation}{ML I}
\begin{multicols}{2}
\subsection*{z-Standardisierung}
$$z_{i}= \frac{x_{i}-\bar{x}}{{s}}$$
$$\bar{z}=0$$
$$s_z=1$$

\subsection*{Zentrierung}

$c_{i} = x_i - \bar{x}$

$\bar{c}=0$

\subsection*{IQ-Standardisierung}
$\mathrm{IQ_i} = z_i \cdot 15 + 100$\
\end{multicols}
\end{formulaBox}
